\section{Introduction}
\label{sec:introduction}

Multimedia channels, e.g., cinema and television (TV), present stories through video, text, image, and audio. By watching how different characters interact with each other in various scenarios, audiences are able to broaden their horizons and experience colourful emotions, such as love, sorrow, and excitement.

As these plot-rich stories (e.g., movies and TV dramas) reflect our real world, they have long been used as devices for the development and assessment of human intelligence (especially for children) in the Cognitive Science community~\citep{dev-theory-mind,child-theory-mind,gunter2005children,bullet-screen-pyschology}. Analogously, a large spectrum of works employed movies and TV dramas to train and test multimodal algorithms~\citep{lsmdc2015cvpr,yms2018cvpr,movienet2020eccv,mad2022cvpr}. However, the results obtained from current methods are generally unsatisfactory, particularly when handling clips. They often only scratch the surface and fail to capture the plot of the cinema and TV.

\newpage

\begin{table*}[t]
\centering
\scriptsize
\begin{tabular}{lrrrrrrrr}
\hline
\multirow{2}{*}{\textbf{Dataset}} & \textbf{Total length} & \multirow{2}{*}{\textbf{Clip \#}} & \textbf{Avg clip length} & \multirow{2}{*}{\textbf{Sentence \#}} & \textbf{Avg sentence length} & \multirow{2}{*}{\textbf{BSC \#}} & \multirow{2}{*}{\textbf{Language}} \\
 & \textbf{(minute)} &  & \textbf{(second)} &  & \textbf{(word)} &  &  \\
\hline
LSMDC 16 & 8,820 & 128,085 & 4.1 & 128,118 & 9.0 & \ding{55} & \texttt{EN} \\
CMD & 76,185 & 33,976 & 2.4 & 35,681 & 14.5 & \ding{55} & \texttt{EN}\\
MoiveNet & 119,992 & \texttt{UNK} & \texttt{UNK} & \texttt{UNK} & \texttt{UNK} & \ding{55} & \texttt{EN} \\
Pororo & 1,231 & 16,066 & \texttt{UNK} & 43,394 & \texttt{UNK} & \ding{55}  & \texttt{EN} \\
SYMON & 52,103 & \texttt{UNK} & 2.2 & 683,611 & 13.1 & \ding{55} & \texttt{EN}\\
\TVD (Ours) & 39,908 & 24,875 & 63.2 & 24,875 & 36.9 & 26,438,007  & \texttt{ZH}\\
\hline
\end{tabular}
\caption{\label{tab:compare}
Comparisons between \TVD and related datasets. In \TVD, ``sentence'' refers to pieces of plot text. 
}
\vspace{-0.3in}
\end{table*}

\begin{wrapfigure}{r}{0.6\textwidth}
\centering
    \includegraphics[width=0.55\textwidth]{figure/fig_example.pdf}
    \caption{
    Illustrated comparisons between a traditional scene-oriented dataset (left) and \TVD (right).
    }
    \label{fig:example}
\vspace{-0.1in}
\end{wrapfigure}

Despite these prior attempts having harvested many valuable insights, the datasets they used exhibit three major problems. \textbf{First and foremost}, to the best of our knowledge, in \textit{all} these corpora segmentation is made based on scenes or shots (i.e., continuous views filmed by one camera without interruption), as illustrated in Fig.~\ref{fig:example}. While such a setup does help reduce the annotation burden, it inevitably breaks the coherence of a plot (an interrelated sequence of events in a story). It has been verified that children at the age of 5 are able to understand isolated scenes of a TV drama, but they have difficulty putting pieces of a plot together~\citep{y5-scene-plot}. Therefore, traditional datasets with scene-oriented annotations may not be sufficient to study a multimodal framework's capability of processing complex stories.
\textbf{Second}, the text signals (e.g., captions) in existing datasets are more simple descriptions of visual targets than vivid narratives of the storyline, as information such as the mental states of roles and situational contexts of events is seldom mentioned. From the perspective of Psycholinguistics, the lack of these elements indicates a deficiency of empathy and sensory perception~\citep{dev-theory-mind,child-theory-mind}. Therefore, conventional datasets are arguably limited in terms of reflecting the machine's utility in the \textit{human world}. 
\textbf{Third}, shown in Tab.~\ref{tab:compare}, clips in most existing corpora are rather short (no more than five seconds on average). As a result, they are unlikely to cover long-form relationships, thus being unsuitable for the development of more advanced systems~\citep{pretrain32022arxiv}.

To bridge the above research gaps, we construct a new dataset (namely \TVD) using large-scale TV drama resources and accompanying Chinese text.
Concretely speaking, \TVD contains 1,106 complete episodes from 83 TV drama series, with comprehensive metadata available. The total video length is 39,908 minutes.
Based on 24,875 professionally written sentences that provide empathic and context-aware plot descriptions (crawled from the web, see~\cref{sec:construction}), 449 human annotators produced 24,875 clips (each corresponding to a standalone plot that links multiple scenes). With an average duration of over 60 seconds, these plot-oriented clips mitigate the aforesaid issues of scene-oriented counterparts and are suitable to assess whether a tested system can model long-form relationships in stories.
Moreover, in \TVD we introduced Bullet Screen Comments (BSCs), a type of real-time comment generated by the audience while watching videos, to provide additional text signals. \TVD's final BSC battery consists of an astonishing number of 26,438,007 sentences, which is also the largest BSC corpus to date. All involved modalities are made parallel with aligned timestamps.

\begin{figure*}[t]
\centering
\includegraphics[width=1.0\textwidth]{figure/fig_overall.pdf}
\caption{
Data structure of \TVD. Sampled from ``\textit{Friends 2nd}'' (Ep. 1), this illustration covers two standalone plots.
}
\label{fig:dataview}
\vspace{-0.2in}
\end{figure*}

Empirically, to demonstrate how \TVD can be employed to boost multimodal research, motivated by the practice and findings in Cognitive Science, we made large-scale experiments on three tasks: Genre Classification, Plot Retrieval, and Plot Text Generation. Note that each task relates closely to one milestone of human intelligence growth (from naive to mature) (see~\cref{sec:tasks}). By adapting state-of-the-art architectures~\citep{meter2022cvpr}, we implemented a unified framework that can be further fine-tuned to solve all three tasks, which naturally serves as a strong baseline for follow-up studies.
On the one hand, we identified the weaknesses of current multimodal algorithms in the cognitive-inspired tasks of different levels.
On the other hand, through extensive ablation studies, we made a list of previously unnoticed and even counter-intuition observations, which may inspire future machine intelligence research.

To summarise, our contribution is three-fold: \textbf{(1)~Dataset}. We construct \TVD, a large-scale cinema/TV dramatic corpus that is segmented according to independent plots for the first time and has the most informative narrative data. It also introduces, to our knowledge, the first \textit{directly} public BSC collection (with 26 million BSCs inside).
Remarkably, \TVD is the first non-English multimodal corpus in the cinema/TV domain. \textbf{(2)~Model}. We released the first \textit{open-source} pre-trained model and weights for processing multimodal cinema/TV-drama data. It is also the first attempt of applying a unified network architecture to solve different tasks in this field. Built with state-of-the-art architectures, our method can serve as a strong baseline in follow-up research. \textbf{(3)~Tasks and Experiments}. Based on Cognitive Science research, we introduced three representative tasks and conducted extensive empirical tests. The novel insights reported (including counter-intuition findings) can benefit future multimodal studies. 

